{\scriptsize\setlength{\tabcolsep}{3pt}
\begin{longtable}{rccccp{4.4cm}}
\caption{Casos de teste de \texttt{reg\_status\_reset.vwf}: reset assíncrono.}
\label{tab:caso-reg-status-reset}\\
\hline $t$ (ns) & \texttt{nrst} & \texttt{status\_reg\_in} & \texttt{status\_reg\_wr\_ena} & \texttt{status\_reg\_out} & Observação \\ \hline \endfirsthead
\hline $t$ (ns) & \texttt{nrst} & \texttt{status\_reg\_in} & \texttt{status\_reg\_wr\_ena} & \texttt{status\_reg\_out} & Observação \\ \hline \endhead
0--20 & 0 & 00000 & 00000 & 00000 &  \\
20--40 & 1 & 11111 & 11111 & 11111 & escreve 11111 \\
40--65 & 1 & 11111 & 00000 & 11111 & mantém \\
65--80 & 0 & 11111 & 00000 & 00000 & reset assíncrono, fora de qualquer borda \\
80--103 & 0 & 11111 & 11111 & 00000 & ainda em reset; escrita habilitada sobre a borda: não escreve \\
103--120 & 1 & 11111 & 11111 & 11111 & solta o reset \\
120--140 & 1 & 10101 & 11111 & 10101 & escreve 10101 (a borda desse ciclo já escreve) \\
140--172 & 1 & 10101 & 00000 & 10101 & mantém \\
172--175 & 0 & 10101 & 00000 & 00000 & pulso de reset de 3 ns, sem borda \\
175--180 & 1 & 10101 & 00000 & 00000 & fim do pulso \\
180--200 & 1 & 10101 & 00000 & 00000 & tudo em 0 \\
\hline
\end{longtable}}
\vspace{-0.9cm}
\fonte{Autoria própria. Saídas do fim de cada intervalo, conforme o modelo de referência e a simulação do Quartus II.}
